\section{Introduction}
\label{sec:intro}

Receipt key-information extraction (KIE) systems on the canonical
ICDAR-2019 SROIE Task-3 benchmark~\cite{huang2019icdar} report
macro-F1 plateauing in the $0.78$--$0.86$ range across architectural
families.  The conventional explanation is that the architecture is
the bottleneck: end-to-end transformer
generators~\cite{kim2022donut}, layout-aware
encoders~\cite{huang2022layoutlmv3}, and modular detect-then-read
pipelines~\cite{yu2020pick} converge on similar numbers despite
substantial parameter-count differences.

This paper asks a different question.  \emph{Is the plateau a
property of the architecture, or of the training data?}  We hold the
DONUT architecture and training recipe fixed and vary only the
training-set composition.  We compare three regimes:

\begin{itemize}
\item \textbf{SROIE-only.}  DONUT fine-tuned on the SROIE 626-image
  train fold under the original Kim et al.\ recipe.  Single-corpus
  baseline.
\item \textbf{SROIE + CORD-v2.}  Joint training on SROIE plus
  CORD-v2~\cite{park2019cord} (Korean restaurant receipts, 30-field
  schema).  Schema diversity expands the training distribution
  beyond a single vendor type.
\item \textbf{SROIE + CORD-v2 + WildReceipt.}  Joint training adding
  WildReceipt~\cite{sun2021wildreceipt} (in-the-wild receipts in
  varied lighting and skew).  Imaging-condition diversity expands
  the training distribution beyond clean scans.
\end{itemize}

\paragraph{Headline result.}  The full multi-dataset regime reaches
\VAR{donut_f1} macro-F1 on canonical SROIE Task-3, an absolute lift
of $+\VAR{multidataset_delta}$ over the single-corpus baseline at
the same parameter count and the same evaluation protocol.

\paragraph{Why this matters.}  Three of the donor corpora's
contributions are decomposable: schema diversity (CORD-v2's 30-field
labels exercise tax / subtotal / item-line structure that SROIE
under-represents), imaging-condition diversity (WildReceipt's
phone-camera receipts), and target-domain alignment (SROIE itself
keeps the test distribution in-domain).  We quantify each
contribution separately in Sec.~\ref{sec:results}.

\paragraph{Contributions.}  This paper makes three contributions:

\begin{itemize}
\item A controlled study of training-set composition's effect on
  SROIE Task-3 F1, holding architecture and training recipe fixed.
\item Quantification of three orthogonal data-axis effects (schema
  diversity, imaging diversity, target alignment) and their
  additivity.
\item A reproducibility manifest that publishes the exact training-
  fold composition (image-level SHA-256 plus per-image schema
  mappings) so an independent lab can reproduce the multi-dataset
  fold from the public corpora.
\end{itemize}

\paragraph{Empirical headline.}  The full multi-dataset DONUT under
the paper-recipe training reaches \VAR{donut_f1} macro-F1 on
canonical SROIE Task-3 (mean across $n=\VAR{n_trials}$ seeds,
paired-bootstrap CI $[\VAR{donut_bootstrap_ci_lo},\,
\VAR{donut_bootstrap_ci_hi}]$).  The lift is decomposed in
Tab.~\ref{tab:headline} into the three contributing data axes.
